\section{讨论与局限}

\subsection{增益来源与机制}

本文的增益可以分解为三步。第一，离线网络提供了强均衡表示：在同一波形与 Pilot 条件下，它相对 DFE-RLS 低 49.6\%–95.1\% 的 BER，相对 LMMSE-FIR 也保持同向优势。第二，CIR 后验把每帧的 Pilot 与 Data 软判决转化为对 116 维物理状态的修正；NMSE 下降 67\%–75\%，说明后验在跟踪信道。第三，跨帧协方差与过程噪声使修正随帧累积，并在信息增量与过程噪声平衡时饱和。5/10/15 dB 的窗口增长与这一机制一致；0 dB 的早期峰值则说明低 SNR 下信息吸收快，但随后被错误传播与信道波动抵消。

一个自然的替代解释是：Bayes 只是多用了 Pilot。这个解释不成立，因为 Frozen 与 Bayes 使用同一 Pilot 预算、同一 acquisition CIR 与同一 tail 规则；Bayes 的额外信息来自 Data 软判决与跨帧后验，而不是更多 Pilot。第二个替代解释是：Bayes 只是多算了算力。两臂共享同一网络前向，Bayes 只增加约 30 ms 的后验；增益不能由网络容量或额外神经参数解释。第三个替代解释是：确认 seed 被挑选。确认 seed 在运行前冻结，10/10 session 无负向，且开发 seed 全部排除；看到结果后没有换 seed 或调参。归因的关键不是 Bayes 用了更多观测，而是它把观测组织成状态后验。Frozen 与 Bayes 的接收波形、Pilot 预算、checkpoint 与 tail 规则完全相同；若增益来自波形或 Pilot，两臂应同时受益。实际观察到的差异只出现在 CIR 后验臂，且 NMSE 与 BER 同向变化，这支持物理状态解释。

\subsection{与既有工作的关系}

本文与三类工作相关。模型驱动的深度接收机把物理模型嵌入网络，ViterbiNet 与深度展开是代表\cite{shlezinger2020,hershey2014}；本文的离线网络同样采用展开结构，但在线阶段不更新网络，而是更新物理状态。Pilot 条件化与元学习接收机关注“用少量 Pilot 快速适应”\cite{park2021,cohen2022}；本文把适应对象从网络参数换成 CIR 后验，因而避免了在线标签与网络权重失稳。经典估计理论给出了线性高斯状态的最优递推\cite{kay1993}；本文的做法是把 Data 软判决的符号方差写入行噪声，使无标签的判决反馈可以安全地进入 CIR 更新。三者的结合点在于：离线表示负责“怎么均衡”，在线物理状态负责“当前信道是什么”，两者解耦后，在线更新只需要低维、可辨识且可解释的状态量。

与更新网络权重相比，本文的在线自由度只有 116 维 CIR 的均值与协方差，参数量远小于网络；这不是因为“占比小”就自然可辨识，而是因为 CIR 由 Pilot 与接收波形直接观测，且在 116 符号记忆下仍可由 ridge/贝叶斯先验正则化。本文的全部在线更新都作用于物理状态，不能由本文结果外推网络权重的在线行为。

\subsection{局限}

第一，0 dB 的 H3 不通过。0 dB 的增益最大且 10/10 session 为正，但窗口斜率为 $-0.001238$/bin，形状是早期收敛、随后波动/饱和，而不是持续增长。四 SNR 平均不能掩盖这一点。主矩阵中 0 dB 的 session 间方差也最大（设计效应约 90.7），10 个 session 对斜率估计仍偏少；若要把 0 dB 写成“增长”，需要新的预注册判据与更多独立 session，不能用本次数据重测。15 dB 的每帧错误数很少，session 级均值对单个错误更敏感；这也是该点绝对效应小、但仍需报告逐 session 分布的原因。10 个 session 足以检出本文观测到的效应，但不足以稳定估计 0 dB 的斜率；后续若要做更窄的区间，应增加独立 session，而不是增加同一 session 的帧数。

第二，确认范围有限。确认只运行一次，样本为 10 session × 128 帧；15 dB 的效应仅 $+0.000260$，虽然单侧下界为正，但绝对错误数很少，估计对 session 波动敏感。实验基于仿真信道，116 符号时延、1200 s 相干时间与慢相位/CFO 都是建模假设，尚未在实测 EME 数据上验证；本文也没有做 CIR 漂移、强径数或弥散比例的单因素外推。把这些物理参数标记为文献/测量来源、建模假设或机制变量，是后续工作的一部分。

第三，实时性未完成。确认脚本没有记录逐帧延迟与 GPU 显存；成本来自单独的 profile（1 seed × 4 SNR × 8 帧）：模型前向约 400 ms，Pilot/Data 后验合计约 30 ms，tail 独立前向约 400 ms，总耗时约 830 ms；tail 复用与独立前向逐位一致，可降到约 435 ms，约为 102.4 ms 帧预算的 4.2 倍。warm-start 与 AMP 扫描没有免费加速：只有 warm-start 8 满足四 SNR 不退化，warm-start 4 在 0/5 dB 更好但在 10/15 dB 退化，warm-start 2 更差，AMP 在 GTX 1650 上无 BER 变化也无加速。因此当前实现不实时，模型前向是剩余瓶颈。

第四，离线复现链不完整。原始训练命令、源码 revision、Python/PyTorch/CUDA 环境、完整优化器超参数与显式随机种子没有随 checkpoint 保存；本文给出的训练配置是重构，不是严格复现。checkpoint 与原始大日志也不在当前工作树内，需要恢复资产后才能独立重跑主矩阵。

第五，传统对照的确认级证据仍缺。3 seed × 16 帧的公平对照说明 Frozen/Bayes 在四 SNR 上均优于 CFO+DD-Phase LMMSE/DFE，但样本仍是开发规模；SC-FDE-MMSE 未做 CFO/相位补偿，只作弱对照。传统方法的确认级比较需要与本文确认同量级的 session 数。

第六，Data 软判决的后验仍有错误传播。符号方差项能降低不可靠符号的权重，但不能消除低 SNR 下的深错误；0 dB 的绝对 BER 仍为 17.13\%。此外，\texttt{block\_phase} 相位块更新在 10×64、10 个新开发 seed 上相对 soft 后验没有可归因增量（causal 四 SNR 配对区间全跨零，same-frame 仅 5 dB 为 $+3.4\times10^{-4}$），因此本文保留 soft 后验作为主方法。同帧重解码的绝对增益更高，但它与“用本帧更新后的 CIR 重新解码本帧”的时序绑定，与本文的跨帧因果协议不同；本文只报告因果结果。

\subsection{未来方向}

下一步有三项。其一，在实测 EME 或更完整的信道测量上复核 116 符号记忆、相干时间与慢相位假设，并把 CIR 漂移、CFO 与强径数作为单因素外推，而不是静默改变当前协议。其二，把物理状态后验做成可部署的实时模块：优先复用 tail 前向（已证明逐位一致），再对模型前向做蒸馏、剪枝、重训、批处理或更大 GPU 的工程优化；在 102.4 ms 预算下，当前 4.2 倍差距需要约 4 倍以上的加速。其三，扩展 session 数与帧数，专门检验 0 dB 的早期收敛/饱和形状与高 SNR 的小效应；若要做增长判据，应在看到新数据前冻结窗口定义与样本量。任何后续扩展都应沿用本文的数据角色隔离与一次性确认规则：开发集用于设计与校准，确认集只运行一次；若改变信道、Pilot、窗口或超参数，应视为新方法并重新预注册。只有在这种约束下，物理状态增益与网络权重增益才能被分开归因。
